\lecture{4}
\title{Sampling, Aliasing, and Quantization}

\begin{document}

\subsection{Lecture: Sampling and Aliasing}

In practice, we might have to settle with discretized samples of data,
rather than the full continuous signal.

\textbf{Definition.} Let $\Delta$ be our \textbf{sampling interval}
and $f_s = \frac{1}{\Delta}$ be our \textbf{sampling frequency}.
Our data is then the \textbf{discrete-time (DT) signal} $x[n] = f(\Delta n)$,
as opposed to the \textbf{continuous-time (CT) signal} $f(t)$.

\begin{remark}
    Just as CT signals have a fundamental period, so too do DT signals---but their
    fundamental periods must be integers!
\end{remark}

But an insufficiently fine sampling might yield misleading results!
\begin{center}
    \frame{\includegraphics[width=0.45\textwidth]{lecture-04/misleading-sampling}}
\end{center}

\textbf{Definition.} If signals $f_1(t) = \cos(\omega_1t)$ and $f_2(t) = \cos(\omega_2t)$
generate the same discrete samples, we say $\omega_1$ and $\omega_2$ are \textbf{aliases}.

\begin{remark}
    We only consider pure cosine signals because
    it turns out they result in the most possible aliasing.
\end{remark}

So in the picture above, the frequencies of the red and green curves are aliases.

\begin{theorem}{Aliasing Classification}
    Frequencies $\omega_1$ and $\omega_2$ are aliases
    if and only if $\omega_2 \Delta = 2\pi k \pm \omega_1 \Delta$ for some $k \in \mathbb{Z}$.
\end{theorem}

\begin{proof}
    Aliasing occurs if and only if we have
    $\cos(\omega_1 \Delta n) = \cos(\omega_2 \Delta n)$
    for all $n \in \mathbb{Z}$.  Note the following:
    \[ \cos(\omega_1 \Delta n) - \cos (\omega_2 \Delta n) =
    2 \sin\left(\frac{\omega_1 \Delta + \omega_2 \Delta}{2} \cdot n\right)
    \sin\left(\frac{\omega_1 \Delta - \omega_2 \Delta}{2} \cdot n\right). \]
    The above is always zero if and only if either $\omega_1 \Delta + \omega_2 \Delta$
    or $\omega_1 \Delta - \omega_2 \Delta$ is an integer multiple of $2\pi$. ~ $\blacksquare$
\end{proof}

One \textbf{anti-aliasing} tactic is to assume that all frequencies
are below the \textbf{Nyquist} frequency $\omega_m := \frac{\pi}{\Delta} = \frac{1}{2}f_s$.
This way, as long as $\omega_1$ and $\omega_2$ are assumed to be below $\omega_m$,
it is impossible for them to be aliases unless they are equal.

Equivalently, every frequency $\omega$ has a \emph{unique} \textbf{baseband} frequency $\omega' < \omega_m$
such that $\omega$ and $\omega'$ are aliases.

\begin{problem}
    Determine the baseband frequency of $f(t) = \cos(4000 t)$ when $\Delta = 0.001$.
\end{problem}

\begin{solution}
    The frequency is $\omega = 4000$, which gives $\omega \Delta = 4$.
    In order to be a valid baseband frequency $\omega'$, we require $\omega' \Delta < \pi$.
    Thus, we take the alias $\omega'$ such that $\omega' \Delta = 2\pi - \omega \Delta = 2\pi - 4$,
    yielding a baseband frequency of $\omega' = 2000\pi - 4000$.
\end{solution}

\begin{remark}
    Check your intuition: when $\Delta$ is smaller,
    the baseband frequency is closer to the original frequency.
\end{remark}

\subsection{Lecture: Quantization}

Another way we can discretize data is by restricting
the number of bits used to represent each individual sample.
\begin{center}
    \frame{\includegraphics[width=0.45\textwidth]{lecture-04/quantized-flower}}
\end{center}

\subsection{Recitation: The Sampling Theorem}

Here's what it looks like to sample
more than 2, exactly 2, and fewer than 2 samples per cycle.
\begin{center}
    \frame{\includegraphics[width=0.45\textwidth]{lecture-04/more-than-two}}
    \frame{\includegraphics[width=0.45\textwidth]{lecture-04/exactly-two}}
    \frame{\includegraphics[width=0.45\textwidth]{lecture-04/less-than-two}}
\end{center}
We can observe the following:
\begin{itemize}
    \item When we sample more than 2 times,
    the samples represent the frequency pretty well.
    \item When we sample exactly 2 times,
    the samples still capture the frequency---but
    practically nothing else.
    \item When we sample fewer than 2 times,
    the samples capture the wrong frequency---it
    appears $\frac{2}{3}$ times the original.
\end{itemize}
In a sense, it takes at least 2 samples per cycle
to preserve the information in a sinusoid.

\begin{theorem}{Sampling Theorem}
    If $f_{\text{max}}$ is the highest nonzero frequency
    in a CT signal, then a sampling rate $f_s \geq 2 f_{\text{max}}$
    is enough to prevent all aliasing.
\end{theorem}

\begin{remark}
    Here's a trick question.
    \begin{quote}
        \textbf{Question.} Humans cannot hear frequencies above 20 kHz.
        Does this mean a sampling rate of 40 kHz suffices?
    \end{quote}
    The answer is no.  None of the frequencies we can hear
    will alias, but frequencies above 20 kHz that we can't hear
    will alias into frequncies that we can hear, thus
    distorting the recording.
\end{remark}

\subsection{Recitation: Dimensional Analysis}

We've been using the word ``frequency'' pretty carelessly so far.
But the introduction of $\Delta$ and $f_s$ complicates things.

\begin{quote}
    \textbf{Definition.}
    The CT signal $f(t) = \sin(K t)$ has an \textbf{angular frequency}
    of $\omega_0 := K$ and a \textbf{cyclical frequency} of $f_0 := \frac{K}{2\pi}$.
    Angular frequency has units \emph{radians per second},
    and cyclical frequency \emph{cycles per second} (or \emph{Hz}).

    \textbf{Definition.}
    The DT signal $f[n] = \sin(K n)$ has an \textbf{angular frequency}
    of $\Omega_0 := K$ and a \textbf{cyclical frequency} of $F_0 := \frac{K}{2\pi}$.
    Angular frequency has units \emph{radians per sample} (dimensionless!),
    and cyclical frequency \emph{cycles per sample} (dimensionless!).
\end{quote}

Cyclical frequency has become relevant with the introduction of $f_s = \frac{1}{\Delta}$,
the sampling frequency (in \emph{samples per second}).

\begin{problem}
    A $5\pi$ Hz CT signal is sampled at an interval of $\Delta = \pi^2$.
    Determine the angular frequency of the resulting DT signal.
\end{problem}

\begin{solution}
    A $5\pi$ Hz CT signal looks like $f(t) = \sin(10\pi^2 t)$ (e.g.).
    The DT signal then looks like $f[n] = \sin(10 \pi^4 n)$,
    from which we can read off an angular frequency of $\Omega_0 = 10\pi^4$.
\end{solution}

\begin{remark}
    Even though the cyclical frequency of a CT signal is measured in Hz,
    the cyclical frequency of a DT signal is not!
\end{remark}

\subsection{Recitation: Periods of DT Signals}

Furthermore, sinusoidal CT signals always have a fundamental period, but DT signals might not!

\begin{quote}
    \textbf{Definition.}
    The CT signal $f(t) = \sin(\omega_0t)$ has a \textbf{fundamental period} $T$,
    with the relationship $T = \frac{2\pi}{\omega_0}$.

    \textbf{Definition.}
    The DT signal $f[n] = \sin(\Omega_0n)$ has a \textbf{fundamental period} $N$
    only if $\frac{2\pi}{\Omega_0}$ is rational,
    in which case $N$ is the numerator of $\frac{2\pi}{\Omega_0}$ in simplest form.
\end{quote}

\textcolor{red!75}{\emph{Warning!}} ~
For CT signals, the relationship
$T = \frac{2\pi}{\omega_0}$ is true,
but for DT signals, the relationship
$N = \frac{2\pi}{\Omega_0}$ is false!

\begin{problem}
    Suppose an $f_0$ Hz CT signal is sampled at a sampling rate of $f_s = 44 \, 100$ Hz.
    What must be true about $f_0$ for the resulting DT signal to be periodic?
\end{problem}

\begin{solution}
    Simply that $f_0 \in \mathbb{Q}^+$.
    The fundamental period is the numerator of $\frac{f_s}{f_0}$.
    So the fundamental period equals $1$ for all of
    $f_0 \in \{44 \, 100, 441 \, 000\}$,
    is fairly small for $f_0 \in \{200, 420\}$,
    and is fairly large for $f_0 \in \{199, 420.01, 0.0042\}$.
\end{solution}

If a DT signal \emph{does} have a fundamental period,
its angular frequency must look like $\Omega_0 = q\pi$ for some $q \in \mathbb{Q}$.
In that case, notice that $\Omega_0' = q\pi \pm 2k\pi$ is \emph{also}
a valid angular frequency!  (And for cosines, we have $\Omega_0' = -q\pi \pm 2k\pi$ too\dots)

(This does not hold for CT signals because $e^{2\pi i n} = 1$
only holds for $n \in \mathbb{Z}$.)

\begin{problem}
    The DT signals $f[n] = \sin(1.2\pi n)$ and $g[n] = \sin(3.2\pi n)$ are the same,
    and they both have fundamental periods.
    What do $\Omega_0$, $F_0$, $N$, $\frac{2\pi}{\Omega_0}$, and $\frac{1}{F_0}$
    equal for each of them?
\end{problem}

\begin{solution}
    Plain computation---but of course, the results are different.
    \begin{itemize}
        \item The ``reduced'' signal $f[n]$ has
        $(\Omega_0, F_0, N, \frac{2\pi}{\Omega_0}, \frac{1}{F_0}) =
        (1.2\pi, 0.6, 5, \frac{5}{3}, \frac{5}{3})$.
        \item The ``unreduced'' signal $g[n]$ has
        $(\Omega_0, F_0, N, \frac{2\pi}{\Omega_0}, \frac{1}{F_0}) =
        (3.2\pi, 1.6, 5, \frac{5}{8}, \frac{5}{8})$.
    \end{itemize}
    Note that $\frac{2\pi}{\Omega_0} = \frac{1}{F_0}$ is always true, by definition.
    Same goes for $\frac{2\pi}{\omega_0} = \frac{1}{\omega_0}$ for CT signals, too.
\end{solution}

Of course, it's a little silly to say that $g[n] = \sin(3.2 \pi n)$
has an ``angular frequency'' of $\Omega_0 = 3.2\pi$ when
we could more cleanly ``alias'' the frequency to $\Omega_0' = 1.2\pi$
without losing any information.  This inspires the following:

\begin{quote}
    \textbf{Definition.} If a DT signal is periodic and has an angular frequency of $\Omega_0$,
    then its \textbf{fundamental angular frequency} is the result of
    reducing $\Omega_0$ modulo $2\pi$.
\end{quote}

(Technically you could alias $\cos(1.6\pi n)$ to $\cos(0.4\pi n)$,
but for ``fundamental angular frequency'', we don't worry about that.)

\textcolor{red!75}{\emph{Warning!}} ~
Even if $\Omega_0$ is the \emph{fundamental} angular frequency,
it is \emph{still} not guaranteed that $N = \frac{2\pi}{\Omega_0}$.

\end{document}